\documentclass[a4paper,11pt]{book}
%
% The current document is based on a document by Marten Wortel
% who is greatly acknowledged for making his document available
%
\usepackage{amsmath,amsthm,amsfonts,mathrsfs,anysize,fancyhdr,epsfig}

\newcommand{\bb}{\mathbb}
\newcommand{\mb}{\mathbf}
\newcommand{\mc}{\mathcal}

\newtheorem{theorem}{Theorem}
\newtheorem{proposition}{Proposition}
\newtheorem{lemma}{Lemma}
\newtheorem*{corollary}{Corollary}

\theoremstyle{definition}
\newtheorem{definition}{Definition}
\newtheorem{conjecture}{Conjecture}
\newtheorem*{example}{Example}
\newtheorem{algorithm}{Algorithm}


\begin{document}

\begin{titlepage}

\thispagestyle{empty}

{\hspace{10cm}
\begin{minipage}{5cm}
  \includegraphics{tudelft.eps}
\end{minipage}}

\begin{center}

\vspace{1.0cm}

\large
\textbf{\large{Delft University of Technology}} \\
\textbf{\large{Faculty of Electrical Engineering, Mathematics and Computer Science}} \\
\textbf{\large{Delft Institute of Applied Mathematics}}

\vspace{2.0cm}

\framebox[16cm]{\raisebox{0cm}[1.5cm][1.5cm]
 {\textbf{
   \begin{minipage}{13cm}\Large
   \begin{center}
       Considering temporal factors in distributed trust.
   \end{center}
   \end{minipage}}}}

\vspace{2.0cm}

\large A thesis submitted to the\\
Delft Institute of Applied Mathematics\\
in partial fulfillment of the requirements

\vspace{1.0cm}

for the degree\\

\vspace{1.0cm}

\textbf{\large{MASTER OF SCIENCE} \\
in \\
\large{APPLIED MATHEMATICS} \\
\vspace{1cm}
by}

\vspace{1cm}

\textbf{PIM OTTE}\\

\vspace*{\stretch{2}}

\textbf{Delft, the Netherlands \\
August 2016} \\

\vspace{1cm}

Copyright \copyright{} 2016 by Pim Otte. All rights reserved.
\end{center}

\end{titlepage}

\newpage

\thispagestyle{empty}

\quad

\newpage

\thispagestyle{empty}

\hspace{10cm}
\begin{minipage}{5cm}
  \includegraphics{tudelft.eps}
\end{minipage}

\begin{center}
\vspace*{\stretch{1}}

\textbf{\large{MSc THESIS APPLIED MATHEMATICS}}

\vspace{1.5cm}

\textbf{``Considering temporal factors in distributed trust''}

\vspace{1.5cm}

PIM OTTE\\

\vspace{1.5cm}

\textbf{\large{Delft University of Technology}}

\vspace{\stretch{5}}

\end{center}

\begin{center}
\begin{tabular}{ll}
\large{\bf{Daily supervisors}} & \large{\bf{Responsible professor}} \\
$\;$ & \\
\large{Dr.~D.~C.~Gijswijt}  & \large{Prof.\ dr.~ir.~K.~Aardal} \\
$\;$ & \\
\large{Dr.~J.~Pouwelse} & \phantom{niets} \\
\phantom{niets} & \phantom{niets} \\
$\;$ & \\
\large{\textbf{Other thesis committee members\qquad}} & \phantom{niets} \\
$\;$ & \\
\large{Drs.~F.~.I.Ctief } & \large{\dots} \\
$\;$ & \\
\large{\dots} & \large{\dots} \\
$\;$ & \\
\phantom{niets} & \phantom{niets} \\
$\;$ & \\
\large{August 2016} & \large{Delft, the Netherlands}
\end{tabular}
\end{center}
\newpage

\thispagestyle{empty}

\quad

\newpage


\chapter*{\centering \begin{normalsize}Abstract\end{normalsize}}

In this thesis, the problem of estimating trust in a distributed network is considered.

\clearpage
\tableofcontents
\pagenumbering{arabic}
\chapter{Introduction}
    
%Insert general gewauwel here

The internet is more distributed than ever [citation needed].

Distributed systems suffer from a tragedy of the commons.

Defending against sybils is interesting in the sense
that it will also imply that there is no gross unfairness in the system.


\section{Base model}

Traditionally, sybil defense mechanisms are built to exploit some graph,
either based in external trust or interactions. In the case of interactions,
information over time is summarized in a single number. However, this loses
a lot of relevant information.

In order to incorporate temporal information in sybil defense mechanisms,
there will be some assumptions needed about the information is available.
Hence, applications which fit the following model will be explored.

\begin{definition}[Ordered Interaction Model]
    An Ordered Interaction Model $M$ consists of two sets:  
    
    \begin{itemize}
        \item $P$, a set of participants
        \item $I$, a set of interactions
    \end{itemize}

    Participants represent entities that can interact with each other.
    An interaction consists of two participants, one or both performing
    work for each other. 
    
    An interaction is defined by 3 pairs of properties:
    \begin{enumerate}
        \item $p, q \in P$, two participants involved in the interaction.
        \item $n_p, n_q \in \bb{N}$, sequence numbers indicating that this is the $n_{\{p,q\}}$th
            interaction of $p$, $q$ respectively.
        \item $v_p, v_q \in \bb{R}_{\geq0}$, representing the value of $p$'s, respectively $q$'s work for the other.
    \end{enumerate}
\end{definition}

  
This model induces a graph that resembles traditional interaction graphs, but preservers the knowledge
of the order of transactions.

\begin{definition}[Ordered Interaction Graph]
    A graph $G_M = (V, E)$ is defined as the Ordered Interaction Graph derived from an Ordered Interaction Model $M = (P,I)$,
    its structure is as follows:

    \begin{itemize}
        \item $V := \{ v_{(p, n_p)} : \exists i \in I$, such that $i$ is the $n_p$th interaction for participant $p\}$\\
        \item $E := \{ (v_{(p, n_q)}, v_{(q, n_q)} : \exists i \in I$, such that $i$ is the $n_p$th interaction
        for $p$ and the $n_q$th interaction for $q$. $\} \cup \{ (v_{(p, n)}, v_{(p, n+1)}) : n \geq 2, v_{(p, n+1)} \in V\}$
    \end{itemize}
\end{definition}
    


\end{document}
